\documentclass[11pt,a4paper]{article}
\usepackage[T1]{fontenc}
\usepackage[utf8]{inputenc}
\usepackage{lmodern,amsmath,amssymb,amsthm,mathtools,microtype,mathrsfs}
\usepackage[margin=25mm,top=23mm,bottom=25mm]{geometry}
\usepackage{booktabs,array,longtable,enumitem,xcolor,fancyhdr}
\usepackage[colorlinks=true,linkcolor=blue!45!black,urlcolor=blue!45!black,citecolor=blue!45!black]{hyperref}
\hypersetup{pdftitle={Reading the Alpoge--Claude RH and Dirichlet-GRH proof: working version 0.44},pdfauthor={ChatGPT},pdfsubject={Partial human-readable reconstruction with explicit arithmetic inputs},pdfkeywords={Riemann hypothesis, Lean, theta lifts, prime sums, Arakelov slopes}}
\setlength{\parindent}{0pt}
\setlength{\parskip}{5pt plus 1pt}
\setlist{topsep=3pt,itemsep=3pt,leftmargin=*}
\setlength{\headheight}{14pt}
\pagestyle{fancy}\fancyhf{}
\fancyhead[L]{\small RH and Dirichlet GRH: a reconstruction}
\fancyhead[R]{\small Working version 0.44}
\fancyfoot[C]{\thepage}
\renewcommand{\headrulewidth}{0.3pt}
\numberwithin{equation}{section}
\newtheorem{theorem}{Theorem}[section]
\newtheorem{proposition}[theorem]{Proposition}
\newtheorem{lemma}[theorem]{Lemma}
\newtheorem{corollary}[theorem]{Corollary}
\theoremstyle{remark}\newtheorem{remark}[theorem]{Remark}
\newcommand{\Rea}{\operatorname{Re}}
\newcommand{\HS}{\mathrm{HS}}
\newcommand{\tr}{\operatorname{tr}}
\newcommand{\codim}{\operatorname{codim}}
\newcommand{\covol}{\operatorname{covol}}
\newcommand{\OK}{\mathcal O_K}
\newcommand{\src}[2]{\href{https://github.com/CocoSgt/riemann-hypothesis-lean/blob/7cbf7a0c32fca1badf05931879a37c3f8a2b03b4/src/#1}{\texttt{#2}}}
\newcommand{\status}[1]{\textbf{Status:} #1}

\begin{document}
\section{Arithmetic points, cyclotomic strings, and the strict cofinite bound}
\label{sec:hida-decomposition-v45}

This chapter proves the original arithmetic-point extraction theorem and
the complete \texttt{Arithmetic\allowbreak Point\allowbreak Decomposition} field, also called
\texttt{end1}.  The same point and the same supplied representation are
retained throughout.  The strict local estimate needed at the last step
follows from the generic-unitary classification already matched in BCS8
and the spherical calculation in GUA12.  It requires no additional
Ramanujan statement or stronger uniform window theorem.

We retain the original \texttt{HypothesesH'}, the coefficient DVR
\(\mathcal O\), the fixed field isomorphism
\(\iota:\overline{\mathbb Q}_\ell\simeq\mathbb C\), the Hida setting
\(s\), level \(V\), auxiliary datum \(Q\), residual representation,
universal determinant ring \(B\), weight object \(\Lambda\), and
\(R_H:\mathrm{UnitaryHeckeData}(Q,\iota,\Lambda)\).
The original base-change and polarized cuspidal attachment foundations,
and their local and period foundations, have exactly the scope stated
in BCS and GUA.  We use CAP6 for the already matched degree-one
Chebotarev statement.  In particular \(\iota\) is used algebraically;
no continuity of \(\iota\) is assumed.

For a number field \(E\), write \(G_E=\operatorname{Gal}(\Omega/E)\),
\(q_w=\#(\mathcal O_E/w)\), and \(\varepsilon=\mathrm{cycBar}\).
The source convention is
\(\varepsilon(\operatorname{Frob}^{\mathrm{ar}}_w)=q_w\).
The predicate \(\mathrm{FrobEig}(\iota,r,w,S)\) includes triviality
on every inertia group above \(w\) and, for every place prolongation
and every arithmetic Frobenius lift \(g\), the identity
\[
 \iota\det(X-r(g^{-1}))=\prod_{z\in S}(X-z).
 \tag{HED.1}
\]
All products over multisets include multiplicities.

\subsection{Extracting the actual representation at a point of the ideal}

% BEGIN HED1
\paragraph{HED1. Factoring a point through the actual Hecke quotient.}
\label{claim:HED1}
Let \(f:A\twoheadrightarrow T\) be a surjective \(\mathcal O\)-algebra
map and let \(x:A\to k\) be another such map, where \(k\) is a field.
If \(\ker f\subseteq\ker x\), there is a unique
\(\mathcal O\)-algebra map \(y:T\to k\) with \(y f=x\).
\begin{proof}
For \(t=f(a)\), set \(y(t)=x(a)\).  If \(f(a)=f(a')\), then
\(a-a'\in\ker f\subseteq\ker x\), proving well-definedness.
Choose preimages of two elements to verify addition and multiplication;
the identity and the \(\mathcal O\)-algebra relation follow from those
for \(f,x\).  Surjectivity proves uniqueness.  Equivalently, this is
the quotient lift \(A/\ker f\to k\) transported through the actual
isomorphism \(A/\ker f\simeq T\).  No section of \(f\) as an algebra
map is required.
\end{proof}
% END HED1

% BEGIN HED2
\paragraph{HED2. Naturality preserves the full determinant law.}
\label{claim:HED2}
Let \(D\) be the original homogeneous multiplicative determinant
polynomial law over \(A\).  For the maps of HED1,
\[
                  (D\otimes_{A,f}T)\otimes_{T,y}k
                              =D\otimes_{A,x}k.
 \tag{HED.2}
\]
Moreover, for a group map \(p:\Gamma'\to\Gamma\), pullback of the
determinant law of a representation \(r\) is the determinant law of
\(r\circ p\).
\begin{proof}
Evaluate both sides on an arbitrary commutative \(k\)-algebra \(C\)
and an arbitrary finite sum in its group algebra.  Both coefficient
structures from \(A\) to \(C\) are identical, because \(yf=x\).
The base-changed polynomial laws therefore have the same value by
their defining naturality.  This proves equality over every scalar
algebra, including those with nilpotents.  Pullback sends a finite sum
\(\sum a_g g\) to \(\sum a_g p(g)\); its matrix under \(r\) is
\(\sum a_g r(p(g))\), exactly the matrix for \(r\circ p\).
Taking determinants proves the second assertion at the same full-law
strength.
\end{proof}
% END HED2

% BEGIN HED3
\paragraph{HED3. Transferring the entire arithmetic-point output.}
\label{claim:HED3}
Suppose the original Hecke algebra \(T=H.TK(\theta,a,K)\) has its
given \(H.\mathrm{Out}\) record for a determinant \(D_0\otimes f\)
and property \(P(\lambda,r)\).  Let \(j:A\to B\) and
\(f:B\twoheadrightarrow T\) satisfy
\(f(j(b))=\operatorname{algebraMap}_{A,T}(b)\).
If \(x:B\to\overline{\mathbb Q}_\ell\) kills \(\ker f\), and
\(xj\) is arithmetic at an admitted \(\lambda\), then one obtains
\(r:H.Gq\to\operatorname{GL}_n(\overline{\mathbb Q}_\ell)\) with
\[
 \mathrm{Continuous}(r),\quad \mathrm{IsSSRep}(r),\quad
 P(\lambda,r),\quad D_0\otimes_x\overline{\mathbb Q}_\ell=\det r.
 \tag{HED.3}
\]
\begin{proof}
Factor \(x=yf\) by HED1.  For every \(b\in A\),
\[
 y(\operatorname{algebraMap}_{A,T}(b))=y(f(j(b)))=x(j(b)).
\]
Thus the weight character induced by \(y\) is exactly the given
character induced by \(xj\).  Apply the literal \(\mathrm{Out.pt}\)
field to this \(y,\lambda\).  It supplies one representation with
continuity, semisimplicity, \(P\), and the evaluated determinant law.
HED2 identifies that law with the last law in (HED.3).
No good-point condition is needed for these four conclusions.
\end{proof}
% END HED3

% BEGIN HED4
\paragraph{HED4. The literal original \texttt{hecke\_ideal\_end}.}
\label{claim:HED4}
Retain all original data with \(n\geq1\), any allowed \(K\), and
\(R_H.\mathrm{IsJ}(B,K,J)\).  For every \(x:B.\mathrm{Pt}\) and
weight \(h\) satisfying \(J\subseteq\ker x\) and
\(x.\mathrm{Arith}(h)\), there are an admitted \(\lambda\) and one
frame representation \(r:G_0\to\operatorname{GL}_n
(\overline{\mathbb Q}_\ell)\) such that
\[
 \mathrm{Continuous}(r),\quad \mathrm{IsSSRep}(r),\quad
 \mathrm{UnitaryPointProperties}(\bar r,\iota,\lambda,r),\quad
                     x.\mathrm{IsRep}(r\circ V.\mathrm{pr}).
 \tag{HED.4}
\]
\begin{proof}
In the top-ideal branch of \(\mathrm{IsJ}\), the premise sends
\(1\in J\) to zero under \(x\), contradicting \(x(1)=1\) in the
nonzero field.  Hence that branch has no such point.

In the other branch, the actual \(\mathrm{IsJ}\) witness supplies
the original Hecke algebra \(T\), its commutativity and coefficient
conditions, a surjection \(f:B\twoheadrightarrow T\), and
\(J=\ker f\).  It also supplies the original determinant and
ordinary-character comparison and the weight-map relation
\[
 f(\Lambda.\mathrm{toPair}(B.u)(b))
                         =\operatorname{algebraMap}_{\Lambda,T}(b).
\]
The increasing-weight part of \(x.\mathrm{Arith}(h)\), together
with AFS8, gives \(\lambda\) and exactly the arithmetic Hecke
weight required by HED3.  All dominant weights of this original
\(hset\) are admitted.

HEP12 supplies the complete \(\mathrm{Out}\) record on precisely
these data; its automorphic premise is the actual one installed in
UPA17.  Apply HED3 with \(j=\Lambda.\mathrm{toPair}(B.u)\).
Its property \(P\) is the full \(\mathrm{UnitaryPointProperties}\),
so the representation carries all four fields, including all
Teichm\"uller families and the complete automorphic witness.
Its full frame determinant equals \(x.D\).  Pulling this equality
back along \(V.\mathrm{pr}\) by HED2 gives exactly the source
predicate \(x.\mathrm{IsRep}(r\circ V.\mathrm{pr})\).
The same \(r\) occurs in every field of (HED.4).
\end{proof}
% END HED4

% BEGIN HED5
\paragraph{HED5. Pullback to the original absolute Galois group.}
\label{claim:HED5}
The original map \(p=V.\mathrm{pr}:G_L\to G_0\) is surjective.
Consequently a continuous semisimple frame representation \(r\)
pulls back to a continuous semisimple representation \(r\circ p\).
\begin{proof}
Given \(g\in G_0\), lift its underlying element through the
original quotient map \(\mathrm{toGS}\).  The condition defining
\(G_0\) says that the lift has trivial \(\eta\)-character.
The original \(V.h\eta\) identifies this kernel with the image of
\(G_L\).  Choose the corresponding element of \(G_L\); its image
under \(p\) is \(g\), since both have the chosen quotient image.
This proves surjectivity on the actual source groups.

The source map \(p\) is continuous, so the composition is continuous.
A vector subspace is invariant under \(r\circ p\) if and only if
it is invariant under \(r\), by surjectivity.  An invariant
complement furnished by semisimplicity of \(r\) is therefore an
invariant complement for the pullback.  No local semisimplification
or new representation is involved.
\end{proof}
% END HED5

\subsection{Frobenius comparison with all original normalizations}

% BEGIN HED6
\paragraph{HED6. Degree-one places split in the unramified CM extension.}
\label{claim:HED6}
Let \(E/E^+\) be the original CM quadratic extension, unramified at
every finite place.  If a finite place \(w\) of \(E\) has absolute
residue degree one, then \(w\) is split over \(E^+\).
\begin{proof}
Put \(v=w|_{E^+}\).  Multiplicativity of residue degrees gives
\(f(w/v)f(v/p)=1\), hence \(f(w/v)=1\).  Unramifiedness gives
\(e(w/v)=1\).  If there were only one prime above \(v\), the
degree formula \(\sum_{w'\mid v}e(w'/v)f(w'/v)=2\) would give
\(1=2\).  Thus there are two primes above \(v\), which is the
literal splitting condition.  The argument applies to every such
\(w\); its only ramification hypothesis is the original one.
\end{proof}
% END HED6

% BEGIN HED7
\paragraph{HED7. A cofinite degree-one Frobenius comparison is global.}
\label{claim:HED7}
Let \(R:G_E\to\operatorname{GL}_n(\overline{\mathbb Q}_\ell)\)
and a finite family \(r_i:G_E\to\operatorname{GL}_{a_i}
(\overline{\mathbb Q}_\ell)\) be continuous.  Suppose that outside
one finite set, every degree-one place admits multisets \(B_{i,w}\)
such that
\[
 \mathrm{FrobEig}(\iota,R,w,\coprod_i B_{i,w}),\qquad
              \mathrm{FrobEig}(\iota,r_i,w,B_{i,w})\quad\text{for all }i.
\]
Then for every \(g\in G_E\),
\[
                  \det(X-R(g))=\prod_i\det(X-r_i(g)).
 \tag{HED.5}
\]
\begin{proof}
CAP6 puts an arithmetic Frobenius at a degree-one place outside
the finite exceptional set into every prescribed open coset of
\(G_E\).  These Frobenius elements are therefore dense.  At each
one, (HED.1) and the multiset union give (HED.5) at its inverse,
first after applying \(\iota\), then over
\(\overline{\mathbb Q}_\ell\) by injectivity of \(\iota\).

For each coefficient, the characteristic polynomial is a polynomial
in the matrix entries and is continuous.  Coefficients of the finite
product on the right are finite sums of products of those continuous
functions.  Inversion on \(G_E\) is continuous.  Since the coefficient
field is Hausdorff, equality extends from the dense set to every
\(g^{-1}\), and hence to every \(g\).  All continuity is used over
\(\overline{\mathbb Q}_\ell\), before applying the merely algebraic
\(\iota\).  Semisimplicity is not needed for this comparison lemma.
\end{proof}
% END HED7

% BEGIN HED8
\paragraph{HED8. The exact spherical Frobenius law for any real shift.}
\label{claim:HED8}
Let \(w\nmid\ell\), let \(\pi_w\) be an actual smooth irreducible
spherical representation with normalized Satake multiset \(S_w\),
and let \(t\in\mathbb R\).  If the original
\(\mathrm{ParAtT}(\iota,r,w,\mathrm{rc}(\pi_w)(t))\) holds,
with its full quantification over \(\mathrm{IsRecT}\) maps, then
\[
 \mathrm{FrobEig}\left(\iota,r,w,
              \{\alpha q_w^{-t}:\alpha\in S_w\}\right).
 \tag{HED.6}
\]
This is the original \texttt{frobTwT} conclusion, for arbitrary real
\(t\), not only an exponent arising from a particular Speh length.
\begin{proof}
Fix any prolongation \(A\) over \(w\).  The original existence
theorem supplies a permissible local correspondence there.  GUA12
computes its spherical parameter: inertia is trivial, monodromy
is zero, and arithmetic Frobenius has eigenvalues \(\alpha^{-1}\).
The real twist multiplies these by \(q_w^t\).  Its parameter is
unramified and hence good by GUA1, so the conditional equality in
\(\mathrm{ParAtT}\) applies to the actual Weil--Deligne parameter
of \(r\).

For clarity, this equality implies triviality of the original
inertia action, as in GUA14.  Zero monodromy makes that action
trivial on an open subgroup; its image is finite.  Every inertia
matrix is thus semisimple in characteristic zero.  The original
Frobenius-semisimplification witness commutes with it and differs
by a nilpotent matrix.  Simultaneous diagonalization makes this
difference zero, giving the trivial original inertia action.
For general Weil elements the same commuting correction preserves
the characteristic polynomial.  Inverting the arithmetic
eigenvalues \(\alpha^{-1}q_w^t\) gives
\(\alpha q_w^{-t}\).  Trivial inertia makes the identity hold for
every lift, and \(A\) was arbitrary.  This proves the entire
\(\mathrm{FrobEig}\) predicate.
\end{proof}
% END HED8

% BEGIN HED9
\paragraph{HED9. Cyclotomic strings have the source's sign and multiplicities.}
\label{claim:HED9}
At \(w\nmid\ell\), the character \(\varepsilon^{-j}\),
\(j\geq0\), is unramified and has geometric Frobenius value
\(q_w^j\).  Therefore
\[
 \mathrm{FrobEig}(\iota,r,w,B)
 \ \Longrightarrow\
 \mathrm{FrobEig}(\iota,r\otimes\varepsilon^{-j},w,
                                     \{zq_w^j:z\in B\}).
 \tag{HED.7}
\]
Every twist of a continuous representation occurring here is continuous.
\begin{proof}
Inertia away from \(\ell\) fixes every \(\ell\)-power root of
unity; arithmetic Frobenius raises them to their \(q_w\)-th powers.
Passing to the defining inverse limit gives
\(\varepsilon(I_w)=1\) and
\(\varepsilon(\operatorname{Frob}^{\mathrm{ar}}_w)=q_w\).
Thus \(\varepsilon^{-j}(g^{-1})=q_w^j\) for an arithmetic
Frobenius \(g\).

For any matrix \(M\) and nonzero scalar \(u\),
\(\det(X-uM)=u^{\dim M}\det((X/u)-M)\).  If the characteristic
polynomial of \(M\) is \(\prod_{z\in B}(X-z)\), this identity
is \(\det(X-uM)=\prod_{z\in B}(X-uz)\), with every multiplicity
unchanged.  Apply it to the actual Frobenius matrix, and use the
triviality of both inertia actions.  The original cyclotomic
character is continuous by its inverse-limit construction; powers,
inversion and scalar matrix multiplication preserve continuity.
\end{proof}
% END HED9

\subsection{The complete decomposition of an arbitrary supplied representation}

% BEGIN HED10
\paragraph{HED10. One initial Galois representation for each cusp.}
\label{claim:HED10}
Let \(E/E^+\), \(\ell\ne2\), \(n\geq1\), and the actual
automorphic unitary subspace \(\sigma\) be as in the original
\texttt{endShape} theorem.  Choose its BCS18 witnesses: cusps
\(P_i\), ranks \(a_i\geq1\), lengths \(d_i\geq1\), actual local
components \(\pi_{i,w}\), and infinity multisets \(m_{i,\tau}\),
with \(\sum_i a_i d_i=n\).
For each \(i\) there is a continuous representation
\(r_i:G_E\to\operatorname{GL}_{a_i}(\overline{\mathbb Q}_\ell)\)
whose original compatibility record has local parameter
\[
                  \mathrm{rc}(\pi_{i,w})\bigl((d_i-n)/2\bigr).
 \tag{HED.8}
\]
All local and period quantifiers belong to that one \(r_i\).
\begin{proof}
GUA9 supplies an infinity weight at every embedding and applies
the Speh-string identity to the BCS infinity multisets.  In
particular \(m_{i,\tau}\) has distinct entries, and
\(z+(d_i-n)/2\in\mathbb Z\) for every entry \(z\).
These are exactly the hypotheses of GUA8 with the single integer
\(e_i=d_i-n=\mathrm{ex}(n,d_i,0)\).  Apply GUA8 to \(P_i\),
its entire specified local-component family, and all its infinity
multisets.  It supplies \(r_i\) and the full compatibility record
with (HED.8), including continuity.

Make this finite family choice once.  The remaining string members
will be \(r_i\otimes\varepsilon^{-j}\); they are determined by
this initial representation.  Their construction does not require
identifying separately chosen Galois representations for different
half shifts.
\end{proof}
% END HED10

% BEGIN HED11
\paragraph{HED11. Matching the given representation at degree-one places.}
\label{claim:HED11}
Let \(R:G_E\to\operatorname{GL}_n(\overline{\mathbb Q}_\ell)\)
be any continuous representation satisfying the original
\(\mathrm{LocGlobT}\) comparison with \(\sigma\) for every split
place and every actual local component.  For the \(r_i\) of HED10,
outside one finite set and at every degree-one place \(w\), there
are multisets
\[
 B_{ij,w}=\{\alpha q_w^{(n-d_i)/2+j}:\alpha\in S_{i,w}\},
                       \qquad 0\leq j<d_i,
 \tag{HED.9}
\]
such that \(R\) has Frobenius multiset \(\coprod_{i,j}B_{ij,w}\)
and \(r_i\otimes\varepsilon^{-j}\) has multiset \(B_{ij,w}\).
\begin{proof}
GUA13 makes all the finitely many \(\pi_{i,w}\) spherical away
from one finite set.  Add the finitely many places above \(\ell\).
At any remaining degree-one place, HED6 gives splitting over
\(E^+\).  The BCS18 split-component assertion supplies the actual
irreducible normalized product of the specified Speh components.
Apply the assumed local comparison of \(R\) to that component.
GUA6 identifies its local parameter, for every permissible
correspondence, with
\[
 \bigoplus_i\bigoplus_{0\leq j<d_i}
       \mathrm{rc}(\pi_{i,w})\bigl((d_i-2j-n)/2\bigr).
\]
GUA14, whose conclusion applies to any continuous representation
with this parameter, gives the Frobenius multiset in (HED.9) for
\(R\), including trivial original inertia and all Frobenius lifts.

For each initial \(r_i\), HED8 applied to (HED.8) gives
\(\{\alpha q_w^{(n-d_i)/2}:\alpha\in S_{i,w}\}\).
HED9 multiplies this by \(q_w^j\), giving (HED.9) for the
\(j\)-th cyclotomic twist.  The common exponent identity is
\[
 \frac{n-1}{2}+j-\frac{d_i-1}{2}
                          =\frac{n-d_i}{2}+j.
\]
Thus the comparison uses the given \(R\), not a representation
chosen to replace it, and includes every positive Speh length.
\end{proof}
% END HED11

% BEGIN HED12
\paragraph{HED12. The full cyclotomic-string characteristic polynomial.}
\label{claim:HED12}
Under HED10--HED11, for every \(g\in G_E\),
\[
 \det(X-R(g))=
     \prod_i\prod_{j=0}^{d_i-1}
             \det\bigl(X-(r_i\otimes\varepsilon^{-j})(g)\bigr).
 \tag{HED.10}
\]
The ranks and lengths are positive, and every \(r_i\) is continuous.
\begin{proof}
Index the finite family in HED7 by the pairs \((i,j)\),
\(0\leq j<d_i\).  HED9 supplies continuity of every twist;
HED11 supplies exactly its cofinite degree-one Frobenius premise.
Hence HED7 proves (HED.10) at every group element.  Positivity and
the rank identity come from the same BCS witnesses, and continuity
of the initial representations was proved in HED10.  Neither a
semisimplification nor an isomorphism of the given \(R\) is needed
to establish the stated characteristic-polynomial identity.
\end{proof}
% END HED12

% BEGIN HED13
\paragraph{HED13. The strict spherical bound from the existing local classification.}
\label{claim:HED13}
Let \(F\) be a nonarchimedean characteristic-zero local field with
residue cardinality \(q>1\).  If \(\pi\) is smooth, irreducible,
generic, unitarisable and spherical, then every element \(\alpha\)
of its normalized Satake multiset satisfies
\[
                         q^{-1/2}<|\alpha|<q^{1/2}.
 \tag{HED.11}
\]
In particular this applies to every spherical local component of
the actual conjugate-self-dual cusps \(P_i\) above.
\begin{proof}
BCS8 gives the actual normalized induction
\[
          \pi\simeq\delta_1\nu^{t_1}\times\cdots\times
                             \delta_b\nu^{t_b},\qquad |t_s|<1/2,
\]
with each \(\delta_s\) unitary and square-integrable.  Since the
maximal-compact fixed space of \(\pi\) is nonzero, the tensor
formula for that fixed space in BCS14 makes every
\(\delta_s\nu^{t_s}\) spherical.  The real determinant twists
are trivial on the maximal compact, so every \(\delta_s\) is
spherical as well.

Use the original square-integrable-segment theorem, retained in
BCS9, to write each \(\delta_s\) as a segment representation.
The spherical-classification argument in GUA12 proves that a
spherical segment has length one and unramified rank-one cuspidal
support: spherical supercuspidals are unramified characters, and
a longer segment would have the same spherical Hecke character
as a determinant character, contradicting uniqueness of its
Langlands multisegment.  Thus \(\delta_s=\chi_s\) is a unitary
unramified character, and \(|\chi_s(\varpi)|=1\).

The normalized induction calculation in BCS15 identifies the
Satake multiset, with multiplicities, as
\(\{\chi_s(\varpi)q^{-t_s}\}_s\).  Its moduli are \(q^{-t_s}\).
Since \(q>1\) and \(-1/2<t_s<1/2\), both inequalities in
(HED.11) are strict.  BCS1 and BCS5 supply precisely unitarisability
and genericity for the chosen cuspidal local components.
This derives the source's needed local Satake window from its
already matched generic-unitary input; it asserts no uniform
distance of the exponents from the endpoints as the place varies.
\end{proof}
% END HED13

% BEGIN HED14
\paragraph{HED14. Singleton strings give the exact cofinite ratio at all places.}
\label{claim:HED14}
If every \(d_i=1\) in HED12, then for cofinally many finite places
\(w\), every multiset \(b\) with
\(\mathrm{FrobEig}(\iota,R,w,b)\) and every \(z,z'\in b\) satisfy
\[
                             |z|<q_w|z'|.
 \tag{HED.12}
\]
There is no final splitting condition and no extra nonzero-root
premise in this assertion.
\begin{proof}
Outside the finite set consisting of the places above \(\ell\)
and the nonspherical places of the finitely many cusps, HED8 gives
for each \(r_i\) the multiset
\[
                  C_{i,w}=\{\alpha q_w^{(n-1)/2}:
                                              \alpha\in S_{i,w}\}.
\]
These comparisons hold at every such place, whether split or not.
The identity (HED.10) is already global.  When \(d_i=1\), it says
\(\det(X-R(g))=\prod_i\det(X-r_i(g))\) for every \(g\).

Given any \(b\) in the statement, choose a prolongation and an
arithmetic Frobenius at \(w\), whose existence is part of the
original local-place formalism.  Evaluate that global identity at
its inverse and compare the \(\mathrm{FrobEig}\) equations.  It gives
\[
                  \prod_{z\in b}(X-z)
                         =\prod_i\prod_{u\in C_{i,w}}(X-u).
\]
For completeness, if \(z\in b\), evaluation at \(X=z\) makes the
right side zero.  Since \(\mathbb C\) is a field and the product
is finite, \(z\) equals an element of some \(C_{i,w}\).  Thus
every such \(z\) is \(\alpha q_w^{(n-1)/2}\) for some \(i\).
HED13 consequently gives
\[
                    q_w^{(n-2)/2}<|z|<q_w^{n/2}.
 \tag{HED.13}
\]
The same lower bound holds for \(z'\).  As \(q_w>0\),
\[
             |z|<q_w^{n/2}
                    =q_w q_w^{(n-2)/2}<q_w|z'|.
\]
The finite exceptional set was chosen before \(b,z,z'\), so all
their quantifiers have the required order.  The positive lower
bound also proves nonvanishing rather than requiring it as an
additional hypothesis.
\end{proof}
% END HED14

% BEGIN HED15
\paragraph{HED15. The complete end-shape theorem and its original wrapper.}
\label{claim:HED15}
Under the original \(\mathrm{HypothesesH'}\), let \(E/E^+\) be CM
and unramified at every finite place, \(\ell\ne2\), \(n\geq1\),
and \(\sigma\) an actual \(\mathrm{IsAutRepU}\) subspace.
For every continuous semisimple \(R:G_E\to\operatorname{GL}_n
(\overline{\mathbb Q}_\ell)\) satisfying the original
\(\mathrm{LocGlobT}\) comparison with every split local component
of \(\sigma\), there are finitely many positive ranks \(a_i\),
positive lengths \(d_i\), and continuous representations \(r_i\)
with (HED.10).  If every \(d_i=1\), the exact cofinite conclusion
(HED.12) holds for every \(\mathrm{FrobEig}\) witness of this \(R\).
\begin{proof}
BCS18 supplies one complete Speh family for the given \(\sigma\)
under the original base-change foundation.  HED10 uses the original
polarized attachment foundation through GUA8 to construct the
initial representations.  HED11 and HED12 establish the global
characteristic-polynomial identity for the supplied \(R\).
HED13 and HED14 prove the singleton-string clause.  This gives
every field of the displayed source theorem, including positivity,
continuity, every group element, and its strict cofinite conclusion.
Semisimplicity of \(R\) is retained exactly as in the target.
The \texttt{GuWrap} specialization uses these same BCS18 and GUA8
constructions, so it adds no new supplier.  All odd primes,
including \(3\), and all positive ranks and lengths remain allowed.
\end{proof}
% END HED15

% BEGIN HED16
\paragraph{HED16. The full arithmetic-point decomposition on the same ideal.}
\label{claim:HED16}
Retain the original data of HED4, the condition \(V.S'=s.S\), and
the same \(R_H.\mathrm{IsJ}(B,K,J)\).  Then the original field
\[
               \mathrm{ArithmeticPointDecomposition}(B,\iota,J)
\]
holds.
Explicitly, every \(x,h\) with \(J\subseteq\ker x\) and
\(x.\mathrm{Arith}(h)\) has a continuous semisimple
\(\rho:G_L\to\operatorname{GL}_n(\overline{\mathbb Q}_\ell)\)
with \(x.\mathrm{IsRep}(\rho)\), and positive ranks and lengths
with the full identity (HED.10) and the implication (HED.12).
\begin{proof}
HED4 supplies \(\lambda,r\) and all four unitary point properties.
Set \(\rho=r\circ V.\mathrm{pr}\).  HED5 proves its continuity
and semisimplicity, and HED4 already gives the required full
\(x.\mathrm{IsRep}\) equality.  The \(\mathrm{aut}\) field of
those same point properties supplies one actual \(\sigma\), its
\(\mathrm{IsAutRepU}\) property, and its \(\mathrm{LocGlobT}\)
comparison with this very \(\rho\) at every split local component.

The original Hida setting has odd \(\ell\), and UPA1 proves that
\(L/L^+\) is unramified at every finite place.  Hence HED15 applies
to \(\sigma,\rho\).  Use its ranks \(a_i\) as the source's
\(d_i\), and its lengths \(d_i\) as the source's \(m_i\).
The initial representations \(r_i\), their continuity and
positivity, (HED.10), and (HED.12) are exactly the remaining
fields of the original arithmetic-point decomposition.
No good-point assumption or change of ideal is introduced.
The source's original \(V.S'=s.S\) hypothesis is retained;
the extraction and comparison lemmas themselves have the broader
scopes explicitly stated above.
\end{proof}
% END HED16

\subsection{Scope of this completion}

HED4 is the original point-extraction supplier; HED16 completes the
\texttt{end1} field on the chosen original ideal.  The stronger
\texttt{AutOffN} window, the reverse low-rank point theorem, finite
twisting and all-level modules have separate proofs and exact input
contracts.  The local strict interval in HED13 supplies no place-uniform
margin and is not substituted for that stronger window.  The present
chapter changes neither the original coefficient ring nor the choice
order needed when the Hida fields are installed on both rank-four Links.

\end{document}
